\documentclass[aspectratio=169,10pt]{beamer}
\usetheme{Madrid}
\usecolortheme{default}
\usepackage[utf8]{inputenc}
\usepackage[T1]{fontenc}
\usepackage[spanish,es-nodecimaldot]{babel}
\usepackage{lmodern,booktabs,amsmath}
\definecolor{tfmblue}{RGB}{25,76,120}
\setbeamercolor{structure}{fg=tfmblue}
\setbeamertemplate{navigation symbols}{}
\setbeamertemplate{footline}{}
\title{Detección no supervisada de anomalías\\en radiografías de tórax}
\subtitle{Un autoencoder sencillo, reproducible y explicable}
\author{Alejandro Rodríguez}
\institute{Máster Universitario en Sistemas Inteligentes\\Universitat de les Illes Balears}
\date{Julio de 2026}

\begin{document}
\begin{frame}\titlepage\end{frame}

\begin{frame}{Problema y pregunta}
  \begin{block}{Problema}
    No es viable disponer de ejemplos etiquetados de todas las posibles
    alteraciones radiológicas.
  \end{block}
  \vspace{0.3cm}
  \textbf{Pregunta:} ¿puede un modelo entrenado solo con radiografías normales
  separar las normales de las no normales?
  \begin{itemize}
    \item Salida: puntuación de anomalía, no diagnóstico.
    \item Prioridad: arquitectura pequeña y fácil de defender.
  \end{itemize}
\end{frame}

\begin{frame}{Objetivos}
  \begin{enumerate}
    \item Obtener y auditar un dataset público adecuado.
    \item Implementar las tres familias de la propuesta: AE, VAE y GANomaly.
    \item Entrenar pesos exclusivamente con imágenes normales.
    \item Separar entrenamiento, calibración y prueba.
    \item Conseguir una separación útil y estable con una solución mínima.
    \item Entregar código, un checkpoint y una demostración.
  \end{enumerate}
\end{frame}

\begin{frame}{Dataset y protocolo}
  \textbf{Chest-RSNA (BMAD):} 26.684 radiografías, resolución original
  1024 por 1024 píxeles.
  \begin{center}
  \begin{tabular}{lrr}
    \toprule
    Partición & Normales & Anómalas \\
    \midrule
    Entrenamiento & 8.000 & 0 \\
    Validación & 70 & 1.420 \\
    Test & 781 & 16.413 \\
    \bottomrule
  \end{tabular}
  \end{center}
  \begin{itemize}
    \item Un canal, 64 por 64 píxeles, intensidades normalizadas.
    \item Pesos: solo entrenamiento normal.
    \item Signos y umbral: validación. Resultado final: test congelado.
  \end{itemize}
\end{frame}

\begin{frame}{El modelo final cabe en una línea}
  \begin{center}
  \Large
  $1\rightarrow8\rightarrow16\rightarrow32\rightarrow16\rightarrow8\rightarrow1$
  \end{center}
  \vspace{0.3cm}
  \begin{itemize}
    \item Tres convoluciones de bajada y tres de subida.
    \item ReLU y sigmoide final; pérdida MAE.
    \item \textbf{16.281 parámetros}, Adam $10^{-3}$, tres épocas.
    \item Sin atención, ensamble, red preentrenada ni clasificador.
  \end{itemize}
\end{frame}

\begin{frame}{Una puntuación transparente}
  El MAE aislado tenía dirección inversa en este dataset. No añadimos otro
  modelo: calibramos dos señales sencillas en validación.
  \[
    s(x)=0{,}5\,z(\mathrm{MAE}(x))+
         0{,}5\,z(\mathrm{intensidad}_{centro}-\mathrm{intensidad}_{borde})
  \]
  \begin{itemize}
    \item Pesos fijos 50/50: no hay búsqueda de hiperparámetros.
    \item Validación decide el signo y el umbral.
    \item Ninguna anomalía actualiza los pesos del autoencoder.
  \end{itemize}
\end{frame}

\begin{frame}{Resultados en tres semillas}
  \begin{center}
  \small
  \begin{tabular}{lrrr}
    \toprule
    Método & Parámetros & AUROC & Bal. acc. \\
    \midrule
    VAE & 740.993 & 0,4506 & 0,5064 \\
    GANomaly & 959.586 & 0,5636 & 0,5484 \\
    Control gradiente & 0 & 0,5981 & 0,5894 \\
    \textbf{AE final} & \textbf{16.281} &
      \textbf{0,7608} & \textbf{0,6790} \\
    \bottomrule
  \end{tabular}
  \end{center}
  \vspace{0.3cm}
  Semillas AE: 0,7608; 0,7603; 0,7614. La reducida dispersión muestra que el
  resultado es estable.
\end{frame}

\begin{frame}{¿Por qué mejora?}
  \begin{itemize}
    \item El AE pequeño tiene menos capacidad para copiar cualquier entrada.
    \item La calibración no presupone que una anomalía siempre tenga más error.
    \item Centro--borde aporta una señal global complementaria y explicable.
    \item La combinación supera cada señal por separado sin añadir otra red.
    \item Tiene unas 59 veces menos parámetros que GANomaly.
  \end{itemize}
  \begin{alertblock}{Lección}
    Una puntuación adecuada puede importar más que aumentar la arquitectura.
  \end{alertblock}
\end{frame}

\begin{frame}{Demostración y reproducibilidad}
  \begin{block}{Entrenamiento completo}
    \texttt{python run.py}
  \end{block}
  \begin{block}{Una radiografía}
    \texttt{python demo.py ruta/a/imagen.png}
  \end{block}
  La demo carga \texttt{modelo\_autoencoder.pt}, reconstruye la imagen, calcula
  la puntuación y la compara con el umbral congelado. Se entregan configuración,
  resultados CSV, tres semillas y un solo checkpoint final.
\end{frame}

\begin{frame}{Limitaciones}
  \begin{itemize}
    \item ``Anómala'' significa no normal según RSNA/BMAD; no diagnostica una
    patología concreta.
    \item El test está muy desequilibrado: 95,46\% de anomalías.
    \item 64$\times$64 puede eliminar hallazgos pequeños.
    \item Centro--borde puede capturar adquisición o encuadre, no solo anatomía.
    \item El test fue consultado durante el desarrollo: falta validación externa
    sin reajuste para estimar generalización clínica.
  \end{itemize}
\end{frame}

\begin{frame}{Conclusiones}
  \begin{itemize}
    \item La propuesta se cumple: AE, VAE y GANomaly, entrenados con normales.
    \item El sistema final obtiene AUROC 0,7608.
    \item La solución ganadora es también la menor: seis capas y 16.281
    parámetros.
    \item Es estable, reproducible y se explica con una fórmula.
    \item Es un prototipo de investigación, no una herramienta clínica.
  \end{itemize}
  \vfill
  \centering\Large Gracias. ¿Preguntas?
\end{frame}
\end{document}
